\documentclass[12pt]{article}
\usepackage[utf8]{inputenc}
\usepackage[spanish]{babel}
\usepackage{amsmath}
\usepackage{amsfonts}
\usepackage{amssymb}
\usepackage{geometry}
\geometry{letterpaper, margin=2.5cm}

\begin{document}

\begin{center}
    \Large \textbf{Evaluación de Examen 1: Unidad I} \\
    \large Física para Ciencias de la Salud (005-1713) \\
    \vspace{0.5cm}
    \textbf{Estudiante:} Bárbara Marcano \quad \textbf{C.I.:} 31.729.255
    \vspace{0.3cm} \\
    \large \textbf{Calificación Final: 9/10 puntos}
\end{center}

\vspace{0.5cm}

\section*{Análisis de Resultados}

\subsection*{1. Ángulo plano (Conversión a radianes)}
\textbf{Procedimiento y resultado:} Aplica correctamente el factor de conversión formal $\left(\frac{\pi \text{ rad}}{180^\circ}\right)$. Simplifica la fracción paso a paso hasta obtener la expresión exacta e irreductible de $\frac{5\pi}{12}$ rad. Procedimiento limpio y directo. \\
\textbf{Veredicto:} \textbf{Correcto} (2/2 pts).

\subsection*{2. Sistemas de coordenadas (Longitud del hueso)}
\textbf{Procedimiento y resultado:} Utiliza formalmente la ecuación de la distancia entre dos puntos $D = \sqrt{(x_2-x_1)^2 + (y_2-y_1)^2}$. Sustituye correctamente las coordenadas $(1,2)$ y $(7,10)$, calculando $\sqrt{6^2 + 8^2} = \sqrt{100} = 10$. El gráfico anexo en la segunda página corrobora correctamente esta información. \\
\textbf{Veredicto:} \textbf{Correcto} (2/2 pts).

\subsection*{3. Vectores (Fuerza resultante)}
\textbf{Procedimiento y resultado:} Plantea correctamente la ecuación para la fuerza resultante $\vec{F}_R = \vec{F}_1 + \vec{F}_2$. Agrupa las variables por componentes ($\hat{i}$ y $\hat{j}$) y realiza la suma algebraica sin errores de signo, obteniendo el vector resultante correcto: $30\hat{i} + 60\hat{j}$ N. \\
\textbf{Veredicto:} \textbf{Correcto} (2/2 pts).

\subsection*{4. Potencias de diez (Notación científica y prefijo)}
\textbf{Procedimiento y resultado:} Convierte la medida de manera adecuada a potencias de diez ($12,5 \times 10^{-6}$). Añade el comentario de que el exponente a la $-6$ tiene una equivalencia en el Sistema Internacional, pero omite la respuesta final específica requerida: nombrar el prefijo (micrómetros, $\mu$m). \\
\textbf{Veredicto:} \textbf{Incompleto / Parcialmente correcto} (1/2 pts).

\subsection*{5. Sistemas de unidades (Conversión de densidad)}
\textbf{Procedimiento y resultado:} Aplica magistralmente los factores de conversión en cadena. Coloca las equivalencias exactas tanto para la masa ($1\text{ kg} / 1000\text{ g}$) como para el volumen ($1.000.000\text{ cm}^3 / 1\text{ m}^3$). El desarrollo aritmético $\frac{1.030.000}{1000}$ arroja el resultado perfecto de $1030 \text{ kg/m}^3$. \\
\textbf{Veredicto:} \textbf{Correcto} (2/2 pts).

\vspace{0.5cm}
\section*{Comentarios Generales y Sugerencias}
¡Excelente examen! La estudiante demuestra un dominio muy sólido en el uso de fórmulas matemáticas, álgebra vectorial y factores de conversión en cadena, presentando un desarrollo ordenado y claro.
\begin{itemize}
    \item La única sugerencia es prestar atención a los detalles teóricos de las instrucciones para asegurar respuestas 100\% completas, como lo fue identificar con nombre propio el prefijo del S.I. en la Pregunta 4.
\end{itemize}

\end{document}